%------------------------------------------------------------------------------%
% \chapter{Revisão}
%------------------------------------------------------------------------------%

%------------------------------------------------------------------------------%
\begin{exercicios}{Expressões Algébricas}
    
    % 1
\begin{exercicio}
  Simplifique as expressões abaixo, reduzindo os termos semelhantes.
  \begin{mathlist}[1]
      \item (3x + 2) + (5x - 4)
      \item (2y - 3) - (4y - 5)
      \item (y^3 - 4y^2 + y - 1) - (3y^2 + y - 6)
      \item (-5z + 2x - 6) + 3(z + 4x + 2)
      \item (2a - 5b + 3c) + (6a + 2ab - 3c)
      \item -2(a - 2b - 3ab) - 4(b + 2a - 2ab)
      \item \frac{x - 2}{2} - (2 - x)
      \item \frac{2}{3} (2x - 1) + \frac{4}{3} (2 - x)
      \item \frac{1}{2} (x + 2y - 4) + \frac{1}{3} (3y - x + 9)
      \item \frac{1}{2} (a - 3ab + 2b) + \frac{1}{3} (a - 3b + 4ab)
  \end{mathlist}
  \begin{answers}
      \begin{mathlist}[2]
        \item 8 x - 2
        \item - 2 y + 2
        \item y^{3} - 7 y^{2} + 5
        \item 14 x - 2 z
        \item 2 a b + 8 a - 5 b
        \item 2 a \left(7 b - 5\right)
        \item \frac{3 x}{2} - 3
        \item 2
        \item \frac{x}{6} + 2 y + 1
        \item \frac{a \left(- b + 5\right)}{6}
    \end{mathlist}
  \end{answers}
\end{exercicio}


% 2
\begin{exercicio}
  Expanda as expressões e simplifique-as.
  \begin{mathlist}[2]
      \item \left(\frac{x}{5}\right) \, \left(\frac{2}{3} - 2x\right)
      \item \left(- \frac{x}{2}\right) \, \left(2 - \frac{3x}{4}\right)
      \item (5x - 3) (2x + 4)
      \item (8 - 3x) (x^2 + 6)
      \item -2(1 - x) \left(3 + \frac{x}{2}\right)
      \item (0,7x - 0,2) (4 - 0,6x)
      \item \left(x - \frac{1}{2}\right) \, \left(\frac{1}{3} - x \right)
      \item \left(\frac{x}{2} - 3 \right) \, \left(\frac{5}{4} + 2x \right)
      \item \left(\frac{2x}{3} - \frac{3}{2}\right) \, \left(\frac{3}{4} - \frac{x}{3}\right)
      \item (12x - 5) (12x + 5)
      \item (3x + 4)^2
      \item (x - \sqrt{3})^2
  \end{mathlist}
  \begin{answers}
      \begin{mathlist}[2]
          \item \frac{2 x \left(- 3 x + 1\right)}{15}
          \item \frac{x \left(3 x - 8\right)}{8}
          \item 10 x^{2} + 14 x - 12
          \item - 3 x^{3} + 8 x^{2} - 18 x + 48
          \item x^{2} + 5 x - 6
          \item - 0,42 x^{2} + 2,92 x - 0,8
          \item - x^{2} + \frac{5 x}{6} - \frac{1}{6}
          \item x^{2} - \frac{43 x}{8} - \frac{15}{4}
          \item - \frac{2 x^{2}}{9} + x - \frac{9}{8}
          \item 144 x^{2} - 25
          \item 9 x^{2} + 24 x + 16
          \item x^{2} - 2 \sqrt{3} x + 3
    \end{mathlist}
  \end{answers}
\end{exercicio}


% 3
\begin{exercicio}
  Efetue os produtos abaixo.
  \begin{mathlist}[2]
      \item (x^{-1} + 3)(x + 2)
      \item (3x^2 + 2)(6 - \sqrt{x})
      \item (\sqrt{x} + 9)(\sqrt{x} - 9)
      \item \left(\frac{2}{\sqrt{x}} - 5\right) (\sqrt{x} - 1)
      \item \left(\frac{1}{x} - 1\right) \left(4 + \frac{1}{x^2}\right)
      \item 3x^2 (x^3 - 2x^2 - 4x +5)
      \item -4x^3 (x^2 + 2x - 1)
      \item xy^2 (2x + 3xy + 4y)
      \item (3x + 5) (3x^3 + 4x^2 + 2)
      \item (2 - x^2) (3x^3 + 6x^2 - x)
      \item \left(2x^2 - \frac{1}{2}\right) (x^2 + 3)
      \item (x^3 + 1) (x^4 - 3x^2 + 2)
      \item (3 - 2y + y^2) (2y^2) - 5y + 4
      \item (x - y + 1) (2x - 4y + 6)
      \item (x^2 + 2y) (3x - 2xy - y)
      \item (2x - 1)^3
      \item (x - 3) (x + 3) (x - 2)
      \item (2w - 3) (w - 1) (3w + 2)
      \item (x^2 + 3) (x^2 - 2) (2x^2 - 5)
      \item (a + 2b) (3a - b) (2a + 3b)
      \item (a - b) (a + b) (a^2 + b^2)
  \end{mathlist}
  \begin{answers}
      \begin{mathlist}[1]
          \item 3 x + 7 + \frac{2}{x}
          \item - 3 x^{\frac{5}{2}} - 2 \sqrt{x} + 18 x^{2} + 12
          \item x - 81
          \item - 5 \sqrt{x} + 7 - \frac{2}{\sqrt{x}}
          \item -4 + \frac{4}{x} - \frac{1}{x^{2}} + \frac{1}{x^{3}}
          \item 3 x^{2} \left(x^{3} - 2 x^{2} - 4 x + 5\right)
          \item 4 x^{3} \left(- x^{2} - 2 x + 1\right)
          \item x y^{2} \left(3 x y + 2 x + 4 y\right)
          \item 9 x^{4} + 27 x^{3} + 20 x^{2} + 6 x + 10
          \item x \left(- 3 x^{4} - 6 x^{3} + 7 x^{2} + 12 x - 2\right)
          \item 2 x^{4} + \frac{11 x^{2}}{2} - \frac{3}{2}
          \item x^{7} - 3 x^{5} + x^{4} + 2 x^{3} - 3 x^{2} + 2
          \item 2 y^{4} - 4 y^{3} + 6 y^{2} - 5 y + 4
          \item 2 x^{2} - 6 x y + 8 x + 4 y^{2} - 10 y + 6
          \item - 2 x^{3} y + 3 x^{3} - x^{2} y - 4 x y^{2} + 6 x y - 2 y^{2}
          \item 8 x^{3} - 12 x^{2} + 6 x - 1
          \item x^{3} - 2 x^{2} - 9 x + 18
          \item 6 w^{3} - 11 w^{2} - w + 6
          \item 2 x^{6} - 3 x^{4} - 17 x^{2} + 30
          \item 6 a^{3} + 19 a^{2} b + 11 a b^{2} - 6 b^{3}
          \item a^{4} - b^{4}
    \end{mathlist}
  \end{answers}
\end{exercicio}


    
\begin{exercicio}
  Expanda as expressões abaixo
  \begin{mathlist}[2]
            \item (x + 2)^2
            \item (3x + 8)^2
            \item (x^2 - \sqrt{5})^2
            \item (2u + 7v)^2
            \item (4 - y)^2
            \item (3 - 2y)^2
            \item (-2 - x)^2
            \item \left(\frac{x}{2} + 2\right)^2
            \item \left(\sqrt{2x} + 1\right)^2
            \item \left(3 - \frac{5}{x}\right)^2
            \item \left(2x - \frac{1}{x}\right)^2
            \item (4 - x^2)^2
            \item (x^2 - x)^2
            \item (2x^2 - y)^2
            \item (x^2 + \sqrt{x})^2
            \item (x - 2)^2 (3 - x)^2
            \item \left(\frac{x + 3}{1 - x}\right)^2
            \item (2x + 1)^3
            \item (3 - y)^3
            \item (2 \sqrt[3]{x} - 3)^3        
  \end{mathlist}
  \begin{answers}
      \begin{mathlist}[2]
        \item x^{2} + 4 x + 4
        \item 9 x^{2} + 48 x + 64
        \item x^{4} - 2 \sqrt{5} x^{2} + 5
        \item 4 u^{2} + 28 u v + 49 v^{2}
        \item y^{2} - 8 y + 16
        \item 4 y^{2} - 12 y + 9
        \item x^{2} + 4 x + 4
        \item \frac{x^{2}}{4} + 2 x + 4
        \item 2 \sqrt{2x} + 2 x + 1
        \item 9 - \frac{30}{x} + \frac{25}{x^{2}}
        \item 4 x^{2} - 4 + \frac{1}{x^{2}}
        \item x^{4} - 8 x^{2} + 16
        \item x^{4} - 2 x^{3} + x^{2}
        \item 4 x^{4} - 4 x^{2} y + y^{2}
        \item 2 x^{\frac{5}{2}} + x^{4} + x
        \item x^{4} - 10 x^{3} + 37 x^{2} - 60 x + 36
        \item \frac{x^{2} + 6x + 9}{x^{2} - 2 x + 1}
        \item 8 x^{3} + 12 x^{2} + 6 x + 1
        \item - y^{3} + 9 y^{2} - 27 y + 27
        \item - 36 x^{\frac{2}{3}} + 54 \sqrt[3]{x} + 8 x - 27
    \end{mathlist}
  \end{answers}
\end{exercicio}

    
\begin{exercicio}
  Efetue os produtos abaixo.
  \begin{mathlist}[2]
    \item \left(x + 4\right) \left(x - 4\right)
    \item \left(5 x + 6\right) \left(5 x - 6\right)
    \item \left(2 x + 7 y\right) \left(2 x - 7 y\right)
    \item \left(- x + 2\right) \left(x + 2\right)
    \item \left(\frac{3 x}{2} - \frac{1}{3}\right) \left(\frac{3 x}{2} + \frac{1}{3}\right)
    \item \left(x - \frac{1}{x}\right) \left(x + \frac{1}{x}\right)
    \item \left(y^{2} + 4\right) \left(y^{2} - 4\right)
    \item \left(z - \sqrt{3}\right) \left(z + \sqrt{3}\right)
    \item \left(\sqrt{x} + 5\right) \left(\sqrt{x} - 5\right)
    \item \!\!\left(2 \sqrt{x} \!-\! \sqrt{5}\right)\!\! \left(2 \sqrt{x}\! +\! \sqrt{5}\right)
  \end{mathlist}
  \begin{answers}
      \begin{mathlist}[2]
        \item x^{2} - 16
        \item 25 x^{2} - 36
        \item 4 x^{2} - 49 y^{2}
        \item - x^{2} + 4
        \item \frac{9 x^{2}}{4} - \frac{1}{9}
        \item \frac{x^{4} - 1}{x^{2}}
        \item y^{4} - 16
        \item z^{2} - 3
        \item x - 25
        \item 4 x - 5
    \end{mathlist}
  \end{answers}
\end{exercicio}


               
    \begin{exercicio}
        O número áureo é uma constante real irracional, definida como a raiz positiva da equação quadrática obtida a partir de
        \[ \frac{x + 1}{x} = x \]
        Reescreva a equação acima como uma equação quadrática e determine o número áureo.
    \end{exercicio}

    \begin{exercicio}
        Calcule as expressões abaixo, simplificando o resultado. Sempre que necessário, 
	suponha que os termos no interior das raízes são não negativos e que os denominadores 
	são diferentes de zero.
            \begin{mathlist}
                \item \frac{2 (x - \sqrt{3}) (x + \sqrt{3})}{5 (x^2 - 3)}
                \item \left(\frac{6}{\sqrt{x}} - \frac{4}{\sqrt{y}}\right) \left(\frac{6}{\sqrt{x}} + \frac{4}{\sqrt{y}}\right) \left(\frac{xy}{4}\right)
                \item \frac{(2 - \sqrt{x}) (\sqrt{x} + 2)}{4 - (x - 2)^2}
                \item \frac{(x - \sqrt{2}) (x + \sqrt{2})}{(x - 1)^2 + 2x - 3}
                \item \frac{(\sqrt{x + 5}) (\sqrt{x} - 5)}{\sqrt{x - 25}}
                \item \frac{(\sqrt{2x} + \sqrt{5}) (\sqrt{2x} - \sqrt{5})}{(x - 5)^2 - x^2}
                \item \frac{(x - 3)^2 - (x - \sqrt{3}) (x + \sqrt{3})}{2 - x}
                \item \frac{(x + 2)^2 - (2 \sqrt{x} - 1) (2 \sqrt{x} + 1)}{(5 + x)^2}
            \end{mathlist}     
    \end{exercicio}
                        
    \begin{exercicio}
        Reescreva as expressões abaixo, colocando algum termo em evidência e simplificando 
    	o resultado sempre que possível. Quando necessário, suponha que os denominadores são não nulos.
        \begin{mathlist}[2]
            \item 4 - 2y
            \item 6x - 3
            \item -4x - 10
            \item 35x - 5z + 15y
            \item -10a + 14ab
            \item x^2 - 2x
            \item 8ab - 12b + 4ab^2
            \item 3x^5 - 9x^4 + 18x^7
            \item \frac{3x}{32} - \frac{21}{4}
            \item \frac{5x}{2} - \frac{x^2}{2}
            \item xy + x^2y^2
            \item 4xy + 8yz - 12w^2y
            \item xy^2 + y^5 + 3zy^3
            \item - \frac{5}{12x} + \frac{10}{3x^3}
            \item (4x - 1)^2 + (4x - 1) 3x
            \item (5x + 1) (x - 2) - 4 (x - 2)
            \item \frac{6x^2 - 24x}{3x}
            \item \frac{x (3 - 2x) - 2(3 - 2x)}{x - 2}
        \end{mathlist}
    \end{exercicio}

    
\begin{exercicio}
  Fatore as expressões.
  \begin{mathlist}[2]
    \item x^{2} - 9
    \item 16 x^{2} - 1
    \item 9 - \frac{x^{2}}{4}
    \item x^{2} - 64 y^{2}
    \item 4 y^{2} - 25
    \item 36 x^{2} - 100
    \item 16 - 49 x^{2}
    \item 9 u^{2} - v^{2}
    \item 25 - x^{8}
    \item x^{4} - x^{2}
    \item \frac{9 x^{2}}{4} - \frac{1}{9}
    \item x^{2} - 16
    \item \frac{36}{y^{2}} - \frac{1}{9}
    \item \left(x - 7\right)^{2} - 4
    \item x^{2} + 10 x + 25
    \item 4 x^{2} - 12 x + 9
    \item 3 x^{2} + 12 x + 12
    \item x^{2} - x + \frac{1}{4}
    \item 16 x^{2} + 40 x y + 25 y^{2}
    \item x^{2} y^{2} - 2 x y + 1
    \item x^{2} - 2 \sqrt{3} x + 3
    \item \frac{x^{2}}{4} + \frac{x}{3} + \frac{1}{9}
  \end{mathlist}
  \begin{answers}
      \begin{mathlist}[2]
        \item \left(x - 3\right) \left(x + 3\right)
        \item \left(4 x - 1\right) \left(4 x + 1\right)
        \item - \frac{\left(x - 6\right) \left(x + 6\right)}{4}
        \item \left(x - 8 y\right) \left(x + 8 y\right)
        \item \left(2 y - 5\right) \left(2 y + 5\right)
        \item 4 \left(3 x - 5\right) \left(3 x + 5\right)
        \item - \left(7 x - 4\right) \left(7 x + 4\right)
        \item \left(3 u - v\right) \left(3 u + v\right)
        \item - \left(x^{4} - 5\right) \left(x^{4} + 5\right)
        \item x^{2} \left(x - 1\right) \left(x + 1\right)
        \item \frac{\left(9 x - 2\right) \left(9 x + 2\right)}{36}
        \item \left(x - 4\right) \left(x + 4\right)
        \item - \frac{\left(y - 18\right) \left(y + 18\right)}{9 y^{2}}
        \item \left(x - 9\right) \left(x - 5\right)
        \item \left(x + 5\right)^{2}
        \item \left(2 x - 3\right)^{2}
        \item 3 \left(x + 2\right)^{2}
        \item \frac{\left(2 x - 1\right)^{2}}{4}
        \item \left(4 x + 5 y\right)^{2}
        \item \left(x y - 1\right)^{2}
        \item \left(x-\sqrt{3}\right)^2
	\item \frac{\left(3 x + 2\right)^{2}}{36}
    \end{mathlist}
  \end{answers}
\end{exercicio}

    
\begin{exercicio}
  Resolva as equações abaixo.
  \begin{mathlist}[1]
        \item 4 \left(x - \frac{3}{4}\right) (x - 6) = 0
        \item (x - 9)^2 = 0
        \item (x - 5) (2 - x) = 0
        \item 4x (x + 8) = 0
        \item 8 \left(x + \frac{1}{2}\right) (x - 4) = 0
        \item (5x + 3) (2x + 7) = 0
        \item \left(\frac{x}{4} - \frac{3}{2}\right) \left(3 - \frac{x}{4}\right) = 0
        \item \sqrt{2} \left(x + \sqrt{2}\right) \left(\frac{x}{\sqrt{2}} - 1\right) = 0
  \end{mathlist}
  \begin{answers}
      \begin{mathlist}[2]
        \item x_1 = \sfrac{3}{4} \qquad x_2 = 6
        \item x_1 = 9
        \item x_1 = 2 \qquad x_2 = 5
        \item x_1 = -8 \qquad x_2 = 0
        \item x_1 = - \sfrac{1}{2} \qquad x_2 = 4
        \item x_1 = - \sfrac{7}{2} \qquad x_2 = - \sfrac{3}{5}
        \item x_1 = 6 \qquad x_2 = 12
        \item x_1 = \sqrt{2} \qquad x_2 = - \sqrt{2}
    \end{mathlist}
  \end{answers}
\end{exercicio}

\begin{exercicio}
    Reescreva as equações do exercício anterior na forma $ax^2 + bx + c = 0$.
    \begin{answers}
        \begin{mathlist}[2]
            \item 4 x^{2} - 27 x + 18 = 0 
            \item x^{2} - 18 x + 81 = 0 
            \item - x^{2} + 7 x - 10 = 0 
            \item 4 x^{2} + 32 x = 0 
            \item 8 x^{2} - 28 x - 16 = 0 
            \item 10 x^{2} + 41 x + 21 = 0 
            \item - \frac{x^{2}}{16} + \frac{9 x}{8} - \frac{9}{2} = 0 
            \item x^{2} - 2 = 0 
        \end{mathlist}
    \end{answers}
\end{exercicio}

\begin{exercicio}
    Encontre as soluções reais das equações abaixo, caso existam.
    \begin{mathlist}[2]
        \item x^2 - 10 = 0
        \item 3x^2 - 75 = 0
        \item 4x^2 + 81 = 0
        \item \frac{x^2}{6} - \frac{24}{9} = 0
        \item (x - 2)^2 = 4^2
        \item (2x - 1)^2 - 25 = 0
        \item (x + 3)^2 = \frac{1}{9}
        \item \left(\frac{x}{2} + 1\right)^2 = \frac{9}{4}
    \end{mathlist}
    \begin{answers}
        \begin{mathlist}[2]
            \item x_1 = \sqrt{10} \quad x_2 =\!-\!\sqrt{10}
            \item x_1 = -5 \qquad x_2 = 5 
            \item \) Não há soluções
            \item x_1 = -4 \qquad x_2 = 4 
            \item x_1 = -2 \qquad x_2 = 6 
            \item x_1 = -2 \qquad x_2 = 3 
            \item x_1 = - \sfrac{10}{3} \quad x_2 = -\sfrac{8}{3}
            \item x_1 = -5 \qquad x_2 = 1 
        \end{mathlist}
    \end{answers}
\end{exercicio}

\begin{exercicio}
    Determine as raízes das equações.
    \begin{mathlist}[2]
        \item x^2 - 4x = 0
        \item 5x^2 + x = 0
        \item x^2 = -7x
        \item 2x^2 - 3x = 0
        \item -3x^2 - \frac{x}{2} = 0
        \item \frac{x^2}{3} - \frac{x}{6} = 0
        \item 2x - \sqrt{2}x^2 = 0
        \item \sqrt{3}x - \frac{x^2}{\sqrt{3}} = 0
    \end{mathlist}
    \begin{answers}
        \begin{mathlist}[2]
            \item x_1 = 0 \qquad x_2 = 4 
            \item x_1 = - \sfrac{1}{5} \qquad x_2 = 0 
            \item x_1 = -7 \qquad x_2 = 0 
            \item x_1 = 0 \qquad x_2 = \sfrac{3}{2} 
            \item x_1 = - \sfrac{1}{6} \qquad x_2 = 0 
            \item x_1 = 0 \qquad x_2 = \sfrac{1}{2} 
            \item x_1 = 0 \qquad x_2 = \sqrt{2} 
            \item x_1 = 0 \qquad x_2 = 3 
        \end{mathlist}
    \end{answers}
\end{exercicio}

\begin{exercicio}
    Usando a fórmula de Bhaskara, determine, quando possível, as raízes reais das equações.
    \begin{mathlist}[2]
        \item x^2 - 6x + 8 = 0
        \item x^2 - 2x - 15 = 0
        \item x^2 + 6x + 9 = 0
        \item x^2 + 8x + 12 = 0
        \item 2x^2 + 8x - 10 = 0
        \item x^2 - 6x + 10 = 0
        \item 2x^2 - 7x - 4 = 0
        \item 6x^2 - 5x + 1 = 0 
        \item x^2 - 4x + 13 = 0
        \item 25x^2 - 20x + 4 = 0
        \item x^2 - 2\sqrt{5}x + 5 = 0
        \item 2x^2 - 2\sqrt{2}x - 24 = 0
        \item 3x^2 -0,3x - 0,36 = 0
        \item x^2 - 2,4x + 1,44 = 0
        \item x^2 + 2x + 5 = 0
        \item(x + 8)^2 + 4x = 0
        \item(3 - 4x) (x + 3) = 9
        \item (2 - 3x) (2x - 5) = 4
    \end{mathlist}    
    \begin{answers}
        \begin{mathlist}[2]
            \item x_1 = 2 \qquad x_2 = 4 
            \item x_1 = -3 \qquad x_2 = 5 
            \item x_1 = -3 
            \item x_1 = -6 \qquad x_2 = -2 
            \item x_1 = -5 \qquad x_2 = 1 
            \item \) Não há soluções
            \item x_1 = - \sfrac{1}{2} \qquad x_2 = 4 
            \item x_1 = \sfrac{1}{3} \qquad x_2 = \sfrac{1}{2} 
            \item \) Não há soluções
            \item x_1 = \sfrac{2}{5} 
            \item x_1 = \sqrt{5} 
            \item x_1 = - 2 \sqrt{2} \qquad x_2 = 3 \sqrt{2} 
            \item x_1 = -0,3 \qquad x_2 = 0,4 
            \item \) Não há soluções
            \item \) Não há soluções
            \item x_1 = -16 \qquad x_2 = -4 
            \item x_1 = - \sfrac{9}{4} \qquad x_2 = 0 
            \item x_1 = \sfrac{7}{6} \qquad x_2 = 2 
        \end{mathlist}
    \end{answers}
\end{exercicio}

\begin{exercicio}
    Determine quantas raízes as equações abaixo possuem.
    \begin{mathlist}[2]
        \item 2x^2 + 12x + 18 = 0
        \item x^2 - 3x + 8 = 0
        \item -2x^2 - 5x + 9 = 0
        \item \frac{x^2}{5} - 2x + 20 = 0
        \item -x^2 + 16x - 64 = 0
        \item 3x^2 - 4x + 1 = 0
    \end{mathlist}      
    \begin{answers}
        \begin{mathlist}[2]
            \item n=1 \qquad x_1 = -3 
            \item n=0 
            \item n=2 \quad  x_1 = \frac{-5 \pm\sqrt{97}}{4}
            \item n=0 
            \item n=1 \qquad x_1 = 8 
            \item n=2 \qquad x_1 = \sfrac{1}{3} \quad x_2 = 1 
        \end{mathlist}
    \end{answers}
\end{exercicio}



    \begin{exercicio}
        Determine para que valores de $m$ as equações abaixo possuem ao menos uma raiz.
        \begin{mathlist}[2]
            \item -x^2 - 8x + m = 0
            \item 4x^2 + 12x + m = 0
            \item 5x^2 - 8x + m = 0
            \item mx^2 + 6x - 15 = 0
            \item mx^2 - 5x + 10 = 0
            \item mx^2 - 6x + 9 = 0
        \end{mathlist}
    \end{exercicio}
    
    \begin{exercicio}
        Sabendo que a equação \[4x^2 -(m-3)x + 1 = 0\] possui exatamente uma raiz, $x$, 
        determine os possíveis valores da constante $m$.
    \end{exercicio}

    \begin{exercicio}
        Sabendo que a equação \[mx^2 +(2m+5)x+(m+3) = 0\] não possui raízes reais em $x$, 
        determine os possíveis valores da constante $m$.
    \end{exercicio}
    
    \begin{exercicio}
        A equação $4x^2 -12x + c = 0$ tem duas raízes reais, $x_1$ e $x_2$. 
        Sabendo que $x_2$ é 2 unidades maior que $x_1$, determine essas raízes, bem como o valor de $c$.
    \end{exercicio}
    
    
\begin{exercicio}
  Determine as raízes das equações.
  \begin{mathlist}[2]
    \item 9 x^{4} - 20 x^{2} + 4 = 0
    \item x^{4} + 4 x^{2} - 5 = 0
    \item x^{4} - 8 x^{2} + 16 = 0
    \item x^{4} + 13 x^{2} + 36 = 0
    \item 4 x^{4} - 37 x^{2} + 9 = 0
    \item x^{4} - 24 x^{2} - 25 = 0
  \end{mathlist}
  \begin{answers}
      \begin{mathlist}[1]
        \item x_1 = \sqrt{2} \quad x_2 = - \sqrt{2} \quad x_3 = - \frac{\sqrt{2}}{3} \quad x_4 = \frac{\sqrt{2}}{3} 
        \item x_1 = -1 \quad x_2 = 1 
        \item x_1 = -2 \quad x_2 = 2 
        \item \) Não há soluções
        \item x_1 = -3 \quad x_2 = - \frac{1}{2} \quad x_3 = \frac{1}{2} \quad x_4 = 3 
        \item x_1 = -5 \quad x_2 = 5 
    \end{mathlist}
  \end{answers}
\end{exercicio}


    
    \begin{exercicio}
        Determine o domínio das expressões.
        \begin{mathlist}[2]
            \item \frac{x}{3x - 8}
            \item \frac{y - 12}{16 - y^2}
            \item \frac{\sqrt{3x} + 1}{2x - 15 + x^2}
            \item \frac{2x}{16 + 9x^2}
            \item \sqrt{5x - 4}
            \item \sqrt{35 - 7x}
            \item \sqrt{x^2 - 8}
            \item \sqrt[3]{x - 7}
            \item \frac{\sqrt{2x - 5}}{20 - 8x}
            \item \frac{\sqrt{9 - x^2}}{x - 1}
            \item \frac{x - 6}{\sqrt{x - 5}}
            \item \frac{\sqrt{49 - x^2}}{\sqrt{x}}
            \item \frac{\sqrt{3 - x}}{\sqrt{x - 2}}
        \end{mathlist}
    \end{exercicio}
    
        
\begin{exercicio}
  Simplifique as expressões, fatorando os termos, caso necessário. 
  Suponha sempre que os denominadores são não nulos.
  \begin{mathlist}[2]
      \item \frac{2x - 6}{x - 3}
      \item \frac{2x - 6}{3 - x}
      \item \frac{x^2 - 3x}{4x - 12}
      \item \frac{3y - 12}{6y - 18}
      \item \frac{2x - 4}{3x - 6}
      \item \frac{x^2 - x^3}{x}
      \item \frac{x^2 + x^4}{3x^3}
      \item \frac{x^2y - xy^2}{xy}
      \item \frac{x^2y - xy^2}{x - y}
      \item \frac{x^2 - 9}{x^2 - 3x}
      \item \frac{2x^2 - 50}{x^3 + 5x^2}
      \item \frac{x^2 - 5x + 4}{x - 4}
  \end{mathlist}
  \begin{answers}
      \begin{mathlist}[2]
        \item 2
        \item -2
        \item \frac{x}{4}
        \item \frac{y - 4}{2 \left(y - 3\right)}
        \item \frac{2}{3}
        \item - x \left(x - 1\right)
        \item \frac{x^{2} + 1}{3 x}
        \item x-y
        \item xy
        \item \frac{x + 3}{x}
        \item \frac{2 \left(x - 5\right)}{x^{2}}
        \item x - 1
    \end{mathlist}
  \end{answers}
\end{exercicio}

\begin{exercicio}
  Calcule as expressões abaixo e simplifique o resultado quando possível.
  \begin{mathlist}[2]
      \item \!\!\!\left(\!\frac{x^2 - 4x - 12}{x + 2}\!\right)\!\!\!\left(\!\frac{2x + 1}{x - 6}       \!\right)
      \item \!\!\!\left(\!\frac{2x^2 + 8x}{x - 5}    \!\right)\!\!\!\left(\!\frac{x^2 - 25}{4x^2 + 20x}\!\right)            
      \item \dfrac{\;\dfrac{8}{\;5x\;}\;}{\;\dfrac{4}{\;35x\;}\;}
      \item \dfrac{2 - \dfrac{3}{4}}{\;\dfrac{1}{2x} - \dfrac{1}{3x}\;}
      \item \dfrac{3u^3 v^3}{v^5 u^2} + \dfrac{u^2}{v^2}
      \item \dfrac{2w^3}{y^2} - \dfrac{w^6 y^3}{2w^3 y^5}
      \item \dfrac{2}{x} + \dfrac{4}{5}
      \item \dfrac{2}{5x} - \dfrac{4}{3}
      \item \dfrac{2}{5x - 1} + \dfrac{3}{7}
      \item \dfrac{x + 3}{1 - xd} - 2
      \item \dfrac{5x}{x - 4} + \dfrac{3}{x + 1}
      \item \dfrac{x}{x^2 - 9} - \dfrac{3x - 1}{x - 3}
  \end{mathlist}          
  \begin{answers}
      \begin{mathlist}[2]
        \item 2 x + 1
        \item \frac{x + 4}{2}
        \item 14
        \item \frac{15 x}{2}
        \item \frac{u \left(u + 3\right)}{v^{2}}
        \item \frac{3 w^{3}}{2 y^{2}}
        \item \frac{2 \left(2 x + 5\right)}{5 x}
        \item - \frac{2 \left(10 x - 3\right)}{15 x}
        \item \frac{15 x + 11}{7 \left(5 x - 1\right)}
        \item - \frac{2 d x + x + 1}{d x - 1}
        \item \frac{5 x^{2} + 8 x - 12}{\left(x - 4\right) \left(x + 1\right)}
        \item - \frac{3 x^{2} + 7 x - 3}{\left(x - 3\right) \left(x + 3\right)}
    \end{mathlist}
  \end{answers}
\end{exercicio}

    
    \begin{exercicio}
        Racionalize os denominadores, supondo que $x$ pertença a um domínio adequado.
        \begin{mathlist}[2]
            \item \frac{x}{\sqrt{3x}}
            \item \frac{1}{\sqrt{2 + x}}
            \item \frac{2}{2 - \sqrt{2x}}
            \item \frac{\sqrt{x} + 5}{\sqrt{x} + 4}
            \item \frac{x - 1}{\sqrt{2x - 1} - 1}
            \item \frac{\sqrt{3x}}{\sqrt{x} - \sqrt{3}}
        \end{mathlist}      
    \end{exercicio}
    
    
\begin{exercicio}
  Resolva as equações.
  \begin{mathlist}[2]
      \item \frac{x - 2}{x + 3} = 0
      \item \frac{2x + 5}{x - 1} = 3
      \item \frac{5x - 2}{1 - 3x} = -1
      \item \frac{3 - \sfrac{x}{2}}{3x + 8} = \frac{1}{4}
      \item \frac{3x + 5}{4x - 5} = -3
      \item \frac{4 - \sfrac{x}{2}}{4x + 1} = 0
      \item \frac{x}{x^2 - 3x + 2} = 0
      \item \frac{2x^2}{x + 5} = 5
      \item \frac{x^2}{3x - 2} = 2x - 1
      \item \frac{2 \left(x - \sfrac{5}{6}\right) + 1}{5x - 3} = \frac{2}{3}
      \item \frac{x^2 - 26}{x^2 - 9} = 10
      \item \frac{x^2 + 1}{x^2 - 4} = 6
      \item \frac{x^2 + 3x - 10}{x^2 - 2} = 5
      \item \frac{x^2 - 1}{x^2 - 2x + 1} = 4
      \item \frac{x^2 - 4x + 9}{x^2 - 5x + 6} = 3
      \item \frac{2}{x + 1} - \frac{4}{x - 1} = 0
      \item \frac{4}{x + 1} + \frac{1}{x - 1} = \frac{5}{x^2 - 1}
      \item \frac{3}{x + 1} + \frac{2}{x - 1} = 3
      \item \frac{2}{x - 4} + \frac{5}{x - 2} = 3
      \item \frac{6}{x - 3} + \frac{5}{x - 4} = 2
      \item \frac{x}{x - 2} - \frac{3}{x + 2} = 1
      \item \frac{1}{2x + 1} + \frac{1}{3x - 1} = \frac{2}{5}
      \item \frac{3}{x - 2} - \frac{2}{x + 3} = \frac{1}{x}
      \item \frac{2}{x + 1} - \frac{2}{2x - 3} = \frac{3}{x}
      \item \frac{4}{x - 5} - \frac{2}{x - 2} = \frac{12}{(x - 5) (x - 2)}
      \item \frac{4}{5 - 2x} - \frac{3}{x} = 0
  \end{mathlist}    
  \begin{answers}
      \begin{mathlist}[2]
        \item x_1 = 2 
        \item x_1 = 8 
        \item x_1 = \sfrac{1}{2} 
        \item x_1 = \sfrac{4}{5} 
        \item x_1 = \sfrac{2}{3} 
        \item x_1 = 8 
        \item x_1 = 0 
        \item x_1 = - \sfrac{5}{2} \qquad x_2 = 5 
        \item x_1 = \sfrac{2}{5} \qquad x_2 = 1 
        \item x_1 = 1 
        \item x_1 = - \sfrac{8}{3} \qquad x_2 = \sfrac{8}{3} 
        \item x_1 = \sqrt{5} \qquad x_2 = - \sqrt{5} 
        \item x_1 = 0 \qquad x_2 = \sfrac{3}{4} 
        \item x_1 = \sfrac{5}{3} 
        \item x_1 = 1 \qquad x_2 = \sfrac{9}{2} 
        \item x_1 = -3 
        \item x_1 = \sfrac{8}{5} 
        \item x_1 = - \sfrac{1}{3} \qquad x_2 = 2 
        \item x_1 = 3 \qquad x_2 = \sfrac{16}{3} 
        \item x_1 = \sfrac{7}{2} \qquad x_2 = 9 
        \item x_1 = 10 
        \item x_1 = - \sfrac{1}{12} \qquad x_2 = 2 
        \item x_1 = - \sfrac{1}{2} 
        \item x_1 = - \sfrac{9}{4} \qquad x_2 = 1 
        \item \) Não há soluções
        \item x_1 = \sfrac{3}{2} 
    \end{mathlist}
  \end{answers}
\end{exercicio}

\begin{exercicio}
  Resolva as equações.
  \begin{mathlist}[2]
      \item \frac{3x + 4}{1 - 5x} + \frac{3}{2} = \frac{5}{6}
      \item \sqrt{3x + 4} = 8
      \item \sqrt{x + 1} = 2x - 1
      \item \sqrt{2x + 1} = x - 1
      \item \sqrt{x - 3} + x = 9
      \item \sqrt{4 - x} + 2 = 3x
      \item 4 \sqrt{3x - 1} = \frac{2}{3} - 2x
      \item \sqrt{5 - x^2} = 3 - 2x
      \item \sqrt{8x + 25} - 2 = 3 - 4x
      \item \sqrt{4x + 4} - x + 2 = 0
      \item 2x + \sqrt{10 - 4x} = 1
      \item \sqrt{4x + 5} - x = 2x + 4
      \item \sqrt{2x + 5} - 1 = x
      \item 3 + \sqrt{45 - 6x} = x
      \item \sqrt{2x^2 + 7} = 2x - 1
      \item \sqrt{25 - 3x^2} = -x
      \item 2 \sqrt{9x^2 - 7} = 6
      \item \sqrt{x^2 + 3} + x = 5
      \item \sqrt{4x^2 - 1} + 1 = 2x
      \item \sqrt{9x^2 - 2} - 3x = 4
      \item (x + 2)^{2/3} = 9
      \item (5x - 6)^{3/2} = 8
      \item (4x^2 - 13)^{3/5} = 27
      \item x^{1/2} - 13x^{1/4} + 12 = 0
      \item 3 \sqrt{x} - \sqrt[4]{x} = 2
      \item x^{1/3} - x^{1/6} = 2
  \end{mathlist}
  \begin{answers}
      \begin{mathlist}[2]
        \item x_1 = 14 
        \item x_1 = 20 
        \item x_1 = \sfrac{5}{4} 
        \item x_1 = 4 
        \item x_1 = 7 
        \item x_1 = \sfrac{11}{9} 
        \item x_1 = \sfrac{1}{3} 
        \item x_1 = \sfrac{2}{5} 
        \item x_1 = 0 
        \item x_1 = 8 
        \item x_1 = - \sfrac{3}{2} 
        \item x_1 = - \sfrac{11}{9} \qquad x_2 = -1 
        \item x_1 = 2 
        \item x_1 = 6 
        \item x_1 = 3 
        \item x_1 = - \sfrac{5}{2} 
        \item x_1 = - \sfrac{4}{3} \qquad x_2 = \sfrac{4}{3} 
        \item x_1 = \sfrac{11}{5} 
        \item x_1 = \sfrac{1}{2} 
        \item x_1 = - \sfrac{3}{4} 
        \item x_1 = 25 
        \item \) Não há soluções
        \item \) Não há soluções
        \item x_1 = 1 \qquad x_2 = 20736 
        \item x_1 = 1 
        \item x_1 = 64 
    \end{mathlist}
  \end{answers}
\end{exercicio}


    

\begin{exercicio}
  Efetue as operações indicadas e simplifique cada expressão.
  \begin{mathlist}
      \item (7x^2 - 2x + 5) + (2x^2 + 5x - 4)
      \item (3x^2 + 5xy + 2y) + (4 - 3xy - 2x^2)
      \item x - \{2x - [-x - (1 - x)]\}
      \item \left(\frac{1}{3} - 1 + e\right) - \left(- \frac{1}{3} - 1 + e^{-1}\right)
      \item - \frac{3}{4} y - \frac{1}{4} x + 100 + \frac{1}{2} x + \frac{1}{4} y - 120
      \item 3 \sqrt{8} + 8 - 2 \sqrt{y} + \frac{1}{2} \sqrt{x} - \frac{3}{4} \sqrt{y}    
      \item \frac{8}{9} x^2 + \frac{2}{3} x + \frac{16}{3} x^2 - \frac{16}{3} x - 2x + 2
      \item (x + 8) (x - 2)
      \item (3a - 4b)^2
      \item (x + 2y)^2
      \item (2x + y) (2x - y)
      \item (x^{1/2} + 1) \left(\frac{1}{2} x^{-1/2}\right) - (x^{1/2} - 1) \left(\frac{1}{2} x^{-1/2}\right)
      \item 2 (t + \sqrt{t})^2 - 2t^2
  \end{mathlist}
  \begin{answers}
      \begin{mathlist}[2]
        \item 9 x^{2} + 3 x + 1
        \item x^{2} + 2 x y + 2 y + 4
        \item - x - 1
        \item \frac{2}{3} + e - \frac{1}{e}
        \item \frac{x}{4} - \frac{y}{2} - 20
        \item \frac{\sqrt{x}}{2} - \frac{11 \sqrt{y}}{4} + 8 + 6 \sqrt{2}
        \item \frac{56 x^{2}}{9} - \frac{20 x}{3} + 2
        \item x^{2} + 6 x - 16
        \item 9 a^{2} - 24 a b + 16 b^{2}
        \item x^{2} + 4 x y + 4 y^{2}
        \item 4 x^{2} - y^{2}
        \item \frac{1}{\sqrt{x}}
        \item 4 t^{\frac{3}{2}} + 2 t
    \end{mathlist}
  \end{answers}
\end{exercicio}


\begin{exercicio}
  Coloque em evidência o maior fator comum de cada expressão.
  \begin{mathlist}[2]
      \item 4x^5 - 12x^4 - 6x^3
      \item 4x^2 y^2 z - 2x^5 y^2 + 6x^3 y^2 z^2
      \item 3x^{2/3} - 2x^{1/3}
      \item e^{-x} - xe^{-x}
      \item 2ye^{xy^2} + 2xy^3 e^{xy^2}
      \item 2x^{-5/2} - \frac{3}{2} x^{-3/2}
      \item \frac{1}{2} \left(\frac{2}{3} u^{3/2} - 2u^{1/2}\right)
  \end{mathlist}
  \begin{answers}
      \begin{mathlist}[2]
        \item 2 x^{3} \left(2 x^{2} - 6 x - 3\right)
        \item - 2 x^{2} y^{2} \left(x^{3} - 3 x z^{2} - 2 z\right)
        \item \sqrt[3]{x} \left(3 \sqrt[3]{x} - 2\right)
        \item -e^{-x} \left(x - 1\right) 
        \item 2 ye^{x y^{2}} \left(x y^{2} + 1\right) 
        \item \frac{-3 + \frac{4}{x}}{2 x^{\frac{3}{2}}}
        \item \frac{\sqrt{u} \left(u - 3\right)}{3}
    \end{mathlist}
  \end{answers}
\end{exercicio}


\begin{exercicio}
  Fatore cada expressão.
  \begin{mathlist}[2]
      \item 6ac + 3bc - 4ad - 2bd
      \item 4a^2 - b^2
      \item 12x^2 - 3y^2
      \item 10 - 14x - 12x^2
      \item 3x^2 - 6x - 24
      \item 12x^2 - 2x - 30
      \item (x + y)^2 - 1
      \item 8a^2 - 2ab - 6b^2
      \item x^6 + 125
  \end{mathlist}
  \begin{answers}
      \begin{mathlist}[2]
        \item \left(2 a + b\right) \left(3 c - 2 d\right)
        \item \left(2 a - b\right) \left(2 a + b\right)
        \item 3 \left(2 x - y\right) \left(2 x + y\right)
        \item - 2 \left(2 x - 1\right) \left(3 x + 5\right)
        \item 3 \left(x - 4\right) \left(x + 2\right)
        \item 2 \left(2 x + 3\right) \left(3 x - 5\right)
        \item \left(x + y - 1\right) \left(x + y + 1\right)
        \item 2 \left(a - b\right) \left(4 a + 3 b\right)
        \item \left(x^{2} + 5\right) \left(x^{4} - 5 x^{2} + 25\right)
    \end{mathlist}
  \end{answers}
\end{exercicio}


\begin{exercicio}
  Efetue as operações indicadas e simplifique cada expressão.
  \begin{mathlist}
      \item (x^2 + y^2) x - xy (2y)
      \item 2kr (R - r) - kr^2
      \item 2 (x - 1) (2x + 2)^3 [4(x - 1) + (2x + 2)]
      \item 5x^2 (3x^2 + 1)^4 (6x) + (3x^2 + 1)^5 (2x)
      \item 4 (x - 1)^2 (2x + 2)^3 (2) + (2x + 2)^4 (2) (x - 1)
      \item (x^2 + 2)^2 [5 (x^2 + 2)^2 - 3] (2x)
      \item (x^2 - 4) (x^2 + 4) (2x + 8) - (x^2 + 8x - 4) (4x^3)
  \end{mathlist}
  \begin{answers}
      \begin{mathlist}[1]
        \item x \left(x - y\right) \left(x + y\right)
        \item - k r \left(- 2 R + 3 r\right)
        \item 32 \left(x - 1\right) \left(x + 1\right)^{3} \left(3 x - 1\right)
        \item 2 x \left(3 x^{2} + 1\right)^{4} \left(18 x^{2} + 1\right)
        \item 32 \left(x - 1\right) \left(x + 1\right)^{3} \left(3 x - 1\right)
        \item 2 x \left(x^{2} + 2\right)^{2} \left(5 x^{4} + 20 x^{2} + 17\right)
        \item - 2 \left(x^{5} + 12 x^{4} - 8 x^{3} + 16 x + 64\right)
    \end{mathlist}
  \end{answers}
\end{exercicio}

    
\end{exercicios}

%------------------------------------------------------------------------------%
